\section{Assignment claims and the evidence behind them}
\label{app:evidence}

A description of a reservation rule, a balance test, an assignment first stage,
and a reconstruction of every seat's eligibility answer different questions.
This appendix records those distinctions. It covers the principal studies
identified through the project's Rajasthan/UP literature ledger, the
\texttt{quota\_spending} literature collection, publisher and author records,
and reference searches for reservation, randomization, rotation, and allocation.
The search includes foundational women's-reservation studies and contrasting
rank-based procedures outside Rajasthan and UP. It is a documented starting
set, not a claim that every published assignment statement has been located.
Locations below refer to the specified version; working-paper and journal
pagination are not interchangeable.

\subsection{Foundational lottery and rotation designs}

\paragraph{Chattopadhyay and Duflo.}
\citet{chattopadhyay2004} describe women's reservation as a randomized policy
experiment. The published abstract states that head positions ``have been
randomly reserved for a woman.'' The institutional discussion and balance
comparisons give the assignment claim empirical content. The relevant
comparison spans West Bengal and Rajasthan, but the field sampling units differ:
the Rajasthan hamlet sample should not be treated as an enumeration of GP
allocation pools. The conclusion needed here is narrower than a claim about
all subsequent elections: their account concerns the institutional arrangements
and samples they studied. An audit of later cycles must recover the later
eligible pools and rotation histories.

\paragraph{Beaman and coauthors.}
The abstract of \citet{beaman2009} states that they ``exploit random assignment
of gender quotas for leadership positions on Indian village councils.'' The
study uses prior exposure to women leaders to examine subsequent electoral
outcomes and attitudes. \citet{beaman2012} use the same broad policy setting to
study aspirations and attainment. Both make assignment and exposure history
central to interpretation. Neither publication's abstract alone establishes a
statewide compliance rate or the distribution of deviations across districts.
For the latter question, one needs complete administrative rosters and a
benchmark preserving the actual assignment constraints, rather than treating
an outcome effect or balance result as a compliance statistic.

\paragraph{Bhavnani.}
\citet{bhavnani2009} describe ``randomly chosen seats in local legislatures''
in the published abstract. The application concerns Mumbai municipal elections,
not Rajasthan or UP GP-head offices. It is relevant to the persistence argument
and to the use of quotas as natural experiments, but its assignment account
cannot validate a rural state's allocation mechanism. Municipal wards and GP
headships remain separate institutional universes.

\subsection{Rajasthan caste and UP assignment studies}

\paragraph{Chauchard.}
\citet{chauchard2014} describes rotation through ranked SC population shares
in Rajasthan (author PDF, pp.~7--8 and Table~1). The design uses
matched GPs near local allocation thresholds. The project's prior reading of
online Appendix~B, pp.~5--11, records checks of timing disagreements and
sensitivity to Census vintages, household versus person counts, rounding, and
boundary changes. The supplementary PDF was not successfully re-fetched during
this drafting pass; these appendix details retain that access limitation.
The main paper and the existing ledger support a comparison of assignment
validation strategies, not the assertion that Chauchard ignored implementation
imperfections. Modern Census shares on modern GP boundaries cannot by themselves
refute an election-year allocation account.

\paragraph{Dunning and Nilekani.}
\citet{dunning2013} use state-specific allocation thresholds in Karnataka,
Rajasthan, and Bihar. The December 2012 author version, pp.~20--22 and Table~2,
reports balance checks and discusses variation in implementation across
subdistricts. The account uses population totals in passages concerning
Rajasthan/Bihar, whereas Chauchard describes population shares. This distinction
must be resolved against operative orders and administrative population lists;
it must not be silently standardized by the researcher. Their checks address
comparability in the study samples. The present geographic spread statistics
address another quantity and do not overturn those checks merely by differing
from a homogeneous-process benchmark.

\paragraph{Das, Mukhopadhyay, and Saroy.}
\citet{das2023}, identification discussion and Tables~1--2 and D1, study
Rajasthan's OBC reservation after SC/ST assignments. Their lottery account
applies to the remaining eligible GPs. They examine covariate balance with block
fixed effects and additional composition checks. This is a substantial
administrative sample, so describing all previous studies as small-sample
exercises would be inaccurate. Balance conditional on blocks and geographic
variation in assignment residuals are different estimands. A reconstructed
benchmark must retain the exclusion of seats already allocated to other
categories.

\paragraph{Chaturvedi, Das, and Mahajan.}
\citet{cdm2025}, pp.~760--763 and Figure~1, describe block-specific caste and
women's assignment in UP's 2015 election using SEC population information.
They report almost 98\% prediction of female assignment and examine the
assignment discontinuity and covariate balance. This note audits the geographic
distribution of disagreements with their supplied prediction fields. Those
fields are valuable replication inputs, but their presence does not establish
that this project has independently reconstructed the legal rule. The statewide
assignment universe, near-cutoff outcome sample, and complete-case residual
audit have different denominators.

\subsection{Evidence that ranking and implementation deserve separate treatment}

\paragraph{Gajwani and Zhang.}
\citet{gajwani2015}, published pp.~239--240, explicitly state: ``the assignment
of reservations in Tamil Nadu is not random.'' Their account orders GPs using
caste and female population ratios within blocks, allocates women within caste
categories, and retains reservations across two terms before rotating. Their
main comparison and RD sensitivity should therefore not be described as
unconditional lottery assignment. This is direct evidence that a national
reservation mandate does not imply one nationwide randomization algorithm.
It says nothing by itself about the authorization of different procedures
within Rajasthan or within UP.

\paragraph{Pathak and Macours.}
\citet{pathak2017}, pp.~746--749, describe Andhra Pradesh rotation and ranking
by the female population ratio. Their outcome design compares the timing of
children's exposure to reserved leadership across cohorts. Assignment histories,
not simply current reservation indicators, define the comparison. A legal
rotation schedule can induce dependence across cycles; rejecting independence
is consequently insufficient to establish a departure from the prescribed
procedure. The DOI was verified against Crossref as
\texttt{10.1086/692114}; this corrects a conflicting DOI in an upstream
bibliography without altering the upstream repository.

\paragraph{Narasimhan and Weaver.}
\citet{weaver2024} study GP delimitation, not caste or women's reservation.
The November 2023 working-paper Appendix~OA.2, pp.~80--81, examines district
first stages and constraints associated with neighboring villages' population
viability and contiguity. The working paper also reports that district
noncompliance classifications change between delimitation rounds. This is a
precedent for investigating geographic implementation variation and for testing
concrete alternative explanations. Its particular viability constraints cannot
simply be imported into a reservation model. Likewise, classifications based on
first-stage significance depend on precision and should not become this note's
measure of latent compliance.

\subsection{Related work by the author}

\citet{soodquota}, public author version, section on rotation and its
implications and Appendix~B, studies women's electoral performance following
reservation in Rajasthan and UP. It documents transition dependence and
recognizes that conditioning on a subsequent open seat can select on an outcome
of prior assignment. The present note separates that dependence from geographic
heterogeneity and from verified legal noncompliance. Failure to reject an
independence test is not a certificate of random assignment; a negative lag
coefficient is not an independently established rotation rule.

\citet{green2023} provide a methodological comparison: extending a natural
experiment across earlier and later congresses changes the interpretation of
an apparent daughters effect. The relevance is the discipline of testing
whether a finding generalizes across periods and cohorts. That paper supplies
no evidence about Indian reservation law or administrative discretion.

\section{Legal provisions and historical applicability}
\label{app:law}

The hierarchy relevant to these two states is the Constitution of India, state
legislation, delegated rules, and applicable orders. A state-level mandate does
not automatically authorize districts to choose different allocation procedures.
The distinction between an authorized implementing officer and discretion over
the governing procedure is central to the research question.

\subsection{The constitutional mandate}

Article~243D of \citet{constitution} distinguishes seats within Panchayats
from offices of their chairpersons. Clause~(4), relevant to GP heads, assigns
the manner of chairperson reservation to state legislation and separately
requires rotation. Its final proviso states:
\begin{quote}
``the number of offices reserved under this clause shall be allotted by
rotation to different Panchayats at each level.''
\end{quote}
The clause sets a minimum women's share and population-based SC/ST reservation
requirements at the state and tier level. It does not supply every eligible
pool, ordering convention, or lottery procedure needed to simulate a particular
election. Nor does rotation, standing alone, determine whether repeated
reservation is impossible after category changes, roster changes, or exhaustion
of an eligible list. Those questions require the operative state provisions.
The cited official compilation is dated 11 November 2025; its publication date
is not the enactment date of Article~243D.

\subsection{Rajasthan: prescribed procedure and delegated execution}

Section~16 of the Rajasthan Panchayati Raj Act supplies the statutory context
for reservation of chairperson offices \citep{rajasthanact}. The India Code
compilation of the election rules describes SC/ST ranking, category overlap,
and sarpanch/pradhan allocation in Rules~7--9 \citep{rajasthanrules}.
Rule~7's account uses descending population percentages and a lower eligibility
threshold, with continuation and restarting of rotating lists. Rule~8 addresses
places appearing in both SC and ST lists. Rule~9 refers office allocation to
that procedure and assigns implementation to a government-authorized officer.
Its women's allocation provision specifies ``by draw of lots.''

These are substantive restrictions on implementation. They should be encoded
rather than replaced with a generic assertion of random assignment. The
compilation does not establish which amendments and supplementary orders applied
to every election in this panel. A compliance reconstruction still needs the
complete historical roster, prescribed quota and rounding, population vintage,
previous reservations since the applicable reset, category ordering, and the
actual authority for any exception. A district-specific procedure unsupported
by those sources is not made lawful merely because a district officer used it.

\subsection{Uttar Pradesh: historical rules remain incomplete}

The state Act and the Panchayat Raj reservation rules are distinct from rules
for municipal wards or for other Panchayat tiers \citep{upact,uprules}.
The India Code extract reviewed in the project provides definitions and quota
rounding, but does not establish a complete historical implementation of
Rule~5 for all GP-head elections. Substituting a similarly named municipal or
Kshettra/Zila Panchayat rule would change the institution under study.

The official notification of 17 March 2021 changes the definition of previous
election to refer to the 2015 general election and subsequent general elections
\citep{up2021}. Its existence makes a mechanically inherited 1995 baseline
inappropriate without further verification. The sequence of orders, amendments,
and any supersession must be matched to the final allocation used for the
election. The current supplied-prediction audit concerns 2015 and is not a
reconstruction of the 2021 procedure.

\subsection{What the legal evidence can establish now}

The legal texts establish that reservation is governed by mandates and
procedures, not that every observed assignment complies with them. The statistical
results establish departures from specified reference processes, not that those
processes equal the law. The next evidentiary step is to compare complete
administrative allocation pools with the procedure and inputs actually applicable
to them. Cases should be classified as source disagreement, a resolved historical
input difference, an authorized exception with a cited provision, or an unresolved
departure. Only the last two categories concern implementation after the rule and
inputs have been established; neither category alone identifies political motive.

\section{Source validation and remaining coverage}

Article titles, author lists, years, journals, volumes, pages and DOIs were
checked against Crossref records where available, and publisher or author
records otherwise. Primary PDFs and the existing page-specific evidence ledger
were used to distinguish assignment descriptions from validation exercises.
The bibliography retains the user-authored quota paper as an undated working
paper rather than inventing a publication year. Legal sources retain their
instrument dates and explicit compilation notes.

This pass does not certify an exhaustive review or complete election-year legal
reconstruction. The remaining search includes additional state-specific
applications identified in the upstream literature collection, unavailable
supplementary materials, and historical government and district allocation
orders. Their absence limits claims about the entire literature and about legal
compliance, rather than supplying evidence for discretion. Short quotations
above identify the precise claims under discussion; surrounding paragraphs are
paraphrases, not transcriptions of inaccessible appendices or statutes.
